% Section 4: presolve and postsolve.
\section{Presolve and postsolve}\label{sec:presolve}

Presolve shrinks a model before any engine sees it. The reductions in
\file{nirnay/presolve/presolve.py} are the classical ones of Andersen and Andersen~\cite{andersen1995}
and Gondzio~\cite{gondzio1997}, restricted to those with an \emph{exact primal and dual}
postsolve. Postsolve therefore returns duals $y$ and reduced costs $z$ for the original model as
well as~$x$. \code{nirnay.solve} applies presolve by default to LPs and MILPs solved by the simplex,
the interior-point method and branch-and-bound. It is not applied to QPs or to PDLP, which works on the
original rows.

\subsection{Reductions}
Presolve runs over columns and rows repeatedly until a pass changes nothing (at most 20 passes). Every
reduction pushes a record onto a stack; postsolve pops the stack. With $\bar c$ the current
(possibly modified) cost vector and $[\bar r^l,\bar r^u]$ the current row bounds:

\begin{description}[style=nextline,leftmargin=1.2em,font=\sffamily\bfseries\color{navy}]
\item[Fixed column] ($\ell_j=u_j$). Substitute $x_j=\ell_j$: $\bar r\leftarrow \bar r - a_j\ell_j$,
  $c_0\leftarrow c_0+\bar c_j\ell_j$. Postsolve: $x_j=\ell_j$; $z_j$ follows from the final $y$.
\item[Empty column] (no nonzero in any live row). Set $x_j$ to the bound its cost prefers:
  $\ell_j$ if $\bar c_j>0$, $u_j$ if $\bar c_j<0$. If that bound is infinite the problem is
  unbounded. Postsolve: $z_j=\bar c_j$.
\item[Empty row] Drop it, or declare infeasibility if $0\notin[\bar r^l_i,\bar r^u_i]$. Postsolve: $y_i=0$.
\item[Redundant row] If the activity bounds $L_i=\sum_j \min(a_{ij}\ell_j,a_{ij}u_j)$ and
  $U_i=\sum_j\max(a_{ij}\ell_j,a_{ij}u_j)$ lie inside $[\bar r^l_i,\bar r^u_i]$, the row can never
  bind: drop it, $y_i=0$. If $[L_i,U_i]$ misses the row interval the problem is infeasible.
\item[Singleton row] ($a_{ij}x_j\in[\bar r^l_i,\bar r^u_i]$ with one live entry). The row becomes
  a bound on $x_j$:
  \begin{equation}\label{eq:singleton-bound}
     \ell_j \leftarrow \max\bigl(\ell_j,\; \lambda\bigr),\quad u_j\leftarrow \min\bigl(u_j,\;\mu\bigr),\qquad
     [\lambda,\mu] = \begin{cases} [\bar r^l_i/a_{ij},\ \bar r^u_i/a_{ij}] & a_{ij}>0,\\
                                   [\bar r^u_i/a_{ij},\ \bar r^l_i/a_{ij}] & a_{ij}<0,\end{cases}
  \end{equation}
  rounded inward for an integer column. The record keeps which of the two bounds the row
  supplied, and (see \cref{sec:presolve-bug}) the column's cost $\bar c_j$ and its live rows at this moment.
\item[Free column singleton] on an equality row. If $x_j$ appears only in row $i$, the row is
  $\sum_k a_{ik}x_k = b_i$, and the bounds of $x_j$ cannot bind (it is free, or the row and the
  other columns' bounds imply its bounds), then $x_j$ can be eliminated with the row:
  \begin{equation}\label{eq:free-singleton}
    x_j = \frac{b_i-\sum_{k\ne j}a_{ik}x_k}{a_{ij}},\qquad
    y_i = \frac{\bar c_j}{a_{ij}},\qquad
    \bar c_k \leftarrow \bar c_k - a_{ik}\,y_i\ (k\ne j),\qquad c_0\leftarrow c_0 + b_i\,y_i .
  \end{equation}
  The dual is known at removal time, because $x_j$ is basic (it is not at a bound), so its reduced
  cost $\bar c_j - a_{ij}y_i$ must be zero.
\item[Integer bounds] (MILP only). Round integer bounds inward and tighten them by
  activity-based propagation over the live rows (\cref{sec:propagation}); see
  Achterberg et al.~\cite{achterberg2020} for the MIP presolve reductions still missing. The LP dual is not needed
  for a MILP, so these reductions need no dual postsolve.
\end{description}

\subsection{Postsolve and dual recovery}
Postsolve starts from the reduced solution: $x$ on kept columns, $x_j$ of removed columns as
recorded, $y$ on kept rows and $y_i=0$ on removed rows. It maintains $z=c-A\T y$ over the
\emph{original} model and undoes the stack in reverse order:
\begin{itemize}
  \item a free column singleton restores $x_j$ and sets $y_i=\bar c_j/a_{ij}$ from \eqref{eq:free-singleton};
  \item a singleton row computes the reduced cost of $x_j$ \emph{in the problem as it was when
    the row was removed},
    \begin{equation}\label{eq:zbar}
      \bar z_j = \bar c_j^{\,\text{rec}} - \sum_{k\in\mathcal{R}_j^{\text{rec}}} a_{kj}\,y_k ,
    \end{equation}
    where $\bar c_j^{\,\text{rec}}$ and $\mathcal{R}_j^{\text{rec}}$ (the rows in which $x_j$ was
    still live) were recorded at removal. If the bound this row created is active and $\bar z_j$
    has the sign that bound can absorb, the reduced cost moves to the row:
    \begin{equation}\label{eq:singleton-dual}
      \bigl(x_j=\lambda \text{ and } \bar z_j>0\bigr)\ \text{or}\ \bigl(x_j=\mu\text{ and }\bar z_j<0\bigr)
      \quad\Longrightarrow\quad y_i = \frac{\bar z_j}{a_{ij}},\qquad z_j\leftarrow z_j - a_{ij}y_i .
    \end{equation}
    Otherwise $y_i=0$ and the bound's multiplier stays in~$z_j$.
  \item empty, redundant and fixed items need nothing: $y_i=0$, and $z_j$ is already correct.
\end{itemize}
Each assignment of $y_i$ subtracts $a_{i\cdot}y_i$ from the running $z$, so at the end
$z=c-A\T y$ holds exactly and \eqref{eq:kkt-lp} can be checked on the original model.

\subsection{An order-dependence bug in singleton-row duals}\label{sec:presolve-bug}
\begin{lessonbox}[Bug: singleton-row duals depended on the order of undoing]
\textbf{Symptom.} With presolve switched on, some postsolved solutions had an optimal $x$ but duals
that violated the sign conditions \eqref{eq:kkt-lp} of the original model. The violations sat on
columns that had lost singleton rows, and they depended on the order in which the reductions were
undone.

\textbf{Cause.} The first version decided \eqref{eq:singleton-dual} from the running reduced cost
$z_j = c_j - \sum_k a_{kj}y_k$ at the moment the record was undone. That value is wrong in two ways:
\begin{enumerate}[leftmargin=1.4em]
  \item it uses the original $c_j$, while the reduced problem that removed the row had cost
        $\bar c_j$, already modified by free-column-singleton substitutions
        \eqref{eq:free-singleton} made \emph{earlier};
  \item it sums over all rows of $x_j$, including rows removed \emph{before} this one. Their duals
        are assigned \emph{later} in postsolve (the stack is undone in reverse), so they are still
        zero. Another singleton row on the same column is the typical case.
\end{enumerate}
Both errors depend on the order of reductions, so the answer depended on that order too.

\textbf{Fix.} Record $\bar c_j$ and the live rows of column $j$ (indices and coefficients) when
the row is removed, and evaluate \eqref{eq:zbar} from the record. Every row in
$\mathcal{R}_j^{\text{rec}}$ is either kept (its dual comes from the engine) or removed
\emph{after} this row, and so undone \emph{before} it. Its dual is therefore final when
\eqref{eq:zbar} is evaluated. By induction over the stack, each undo step turns an optimal
primal--dual pair of the smaller problem into one of the larger problem.
\end{lessonbox}

The simplex was changed at the same time to recompute $y=B\mT c_B$ exactly for the true costs
when it returns an optimal basis (\file{lp/simplex.py}, \code{\_result}). Its iteration loop
updates reduced costs incrementally and refreshes $y$ only at refactorisation, and the dual
loop may shift costs (\cref{sec:simplex-safeguards}), so the $y$ held at the end is not exactly the
dual of the final basis for the original costs. Postsolve needs exact duals.

\ifnum\haveDualCheck=1
\paragraph{Check} \file{tools/dual\_check.py} solves every Netlib LP with presolve and the dual
simplex (\SI{60}{s} limit) and tests primal feasibility and the sign conditions
\eqref{eq:kkt-lp} on the original model, relative to $1+\|c\|_\infty$. On the run used for this
report, \nDualOk{} of \nDualTot{} problems passed with \nDualBad{} failures. The largest relative
reduced-cost violation was \dualMaxZ{} and the largest row-dual sign violation \dualMaxY{}. The remaining
\nDualStatus{} problems, \dualStatusNames, hit the \SI{60}{s} limit of the check and were not tested.
\fi

\subsection{What presolve removes}
\ifnum\haveCensus=1
Running presolve alone on the \nCensus{} Netlib problems removes \censusRowsPct\% of all rows,
\censusColsPct\% of all columns and \censusNnzPct\% of all nonzeros (median \censusMedRowsPct\% of
rows per problem, up to \censusMaxRowsPct\% on \code{\censusMaxRowsName}). \censusUntouched{}
problems are not reduced at all. \Cref{tab:presolve-rules} counts the reductions by rule. The whole
census takes \censusTime\,s, with \code{\censusMaxTimeName} the slowest at \censusMaxTime\,s, because
presolve is written in plain Python loops.

\begin{table}[ht]
\centering\small
\caption[Presolve reductions on Netlib]{Presolve reductions on the Netlib LPs, by rule: how many
rows or columns each rule removed in total, and on how many problems it fired at least once.}
\label{tab:presolve-rules}
\begin{tabular}{@{}lrr@{}}
\toprule
Rule & Removed & Problems\\
\midrule
\rows{data/presolve_rules.tex}
\bottomrule
\end{tabular}
\end{table}
\fi

These are only the simplest reductions. The comparator study (\cref{sec:competitors}) shows that
doubleton equations, forcing rows, dominated columns and parallel rows account for most of what
HiGHS's presolve removes. Those are next on the roadmap (\cref{sec:roadmap}).

\ifnum\haveSpxPresolve=1
\paragraph{Effect on the dual simplex}
The simplex sweep with presolve is \file{\presolveFile}\ifnum\presolvePartial=1{}, which was still
running when this report was built, with \presolveDone{} of the \nNetlibFiles{} Netlib problems
finished\fi. It solves \presolveSolved{} of its \presolveDone{} problems (failures: \presolveFailNames). On the
\nPresolveCommon{} problems solved both with and without presolve, presolve reduces the dual
iterations by a factor of \presolveIterRatio{} (geometric mean of the ratios). Timed against HiGHS
inside each sweep, which cancels differences in machine load between the two runs, the shifted-mean
ratio is \presolveRatioHighs{} with presolve and \presolveRatioHighsVtwo{} without, and the geometric
mean of per-instance ratios is \presolveGmrHighsCommon{} with presolve and
\presolveGmrHighsCommonVtwo{} without.
(The raw times are not comparable: HiGHS itself ran \presolveLoad{} times slower, in geometric mean,
during the sweep with presolve, because other sweeps shared the machine.)
\ifnum\presolveFaster=1 Presolve therefore pays for itself on these problems.\fi
\ifnum\presolveFaster=0 Presolve therefore roughly breaks even in shifted geometric mean, which weights
the longer solves\ifnum\presolveSmallWorse=1, and it costs time on the many small problems, where the
per-instance ratio rises\fi. Its passes are pure Python loops, which on small problems cost about as
much as the iterations they save, or more. Compiling presolve is part of the first roadmap item, and
the payoff should come from the stronger reductions still to be written.\fi
\ifnum\presolveFaster=-1 Presolve therefore does not yet pay for itself. Its passes are pure Python
loops, and on these small problems they cost more than the iterations they save. Compiling presolve is
part of the first roadmap item.\fi
\fi
